\uuid{riuM}
\titre{Théorème de Bézout}
\auteur{bodin ; croizat ; sacre}
\organisation{exo7}
\licence{CC BY-NC-SA 4.0}
\datecreate{2022-03-08}
\niveau{L1}
\module{Algèbre}
\chapitre{Arithmétique dans Z}
\sousChapitre{Nombres premiers, nombres premiers entre eux}
\difficulte{2}
\theme{Théorème de Bézout}
\exo7qcmid{hilisit.arithmetique.022}

\contenu{
\begin{question}
\qclassification{103.04}
 Pour $a=41$ et $b=7$, on a notamment l'égalité $a \times (-3) + b \times 18 = 3$. Que peut-on en conclure ?
    \begin{answers}
        \bad{$\pgcd(a,b)=3$.}
        \good{$\pgcd(a,b)$ est un diviseur de $3$.}
        \good{Comme $3$ ne divise pas $7$ alors $a$ et $b$ sont premiers entre eux.}
        \bad{$-3$ et $18$ sont premiers entre eux.}
    \end{answers}
    \begin{explanations}
    Comme $\pgcd(a,b)$ divise $a$ et $b$, il divise aussi $a \times (-3) + b \times 18 = 3$. Donc $\pgcd(a,b)$ est un diviseur de $3$ : c'est donc soit $3$, soit $1$. Mais puisque $3$ ne divise pas $b$ (ni $a$ d'ailleurs), on a donc $\pgcd(a,b)=1$ : $a$ et $b$ sont premiers entre eux.\\
    Enfin, les nombres $-3$ et $18$ sont divisibles par $3$.
    \end{explanations}
\end{question}
}
